% KMP (Knuth-Morris-Pratt) Algorithm

Pattern matching in $O(N + M)$ time complexity using a prefix function (pi array).

\begin{lstlisting}
vector<int> prefix_function(string s) {
    int n = s.length();
    vector<int> pi(n, 0);
    for (int i = 1; i < n; i++) {
        int j = pi[i - 1];
        while (j > 0 && s[i] != s[j]) j = pi[j - 1];
        if (s[i] == s[j]) j++;
        pi[i] = j;
    }
    return pi;
}

vector<int> kmp_search(string text, string pattern) {
    string s = pattern + "#" + text;
    int n = s.length(), m = pattern.length();
    vector<int> pi = prefix_function(s);
    vector<int> matches;
    for (int i = m + 1; i < n; i++) {
        if (pi[i] == m) {
            matches.push_back(i - 2 * m); // 0-indexed position in text
        }
    }
    return matches;
}
\end{lstlisting}

To find all occurrences of a \lstinline|pattern| inside a \lstinline|text|, call \lstinline|kmp_search(text, pattern)|. It returns a vector containing the starting 0-indexed positions of all valid matches.
